\section{The Empirical Frontier, 2010--2025}\label{sec:empirics}
% Every number below is a macro of code/paper1_macros.py. Every qualitative statement (which results hold in all
% four views, which survive a familywise correction) is checked there against notes/empirics_v3.json (block CLAIMS);
% the script stops when a refreshed estimate no longer supports a sentence, and names the sentence.

This section measures the frontier that the model describes and confronts five of its implications: the shape of the frontier, the growth technology, the decline of the attainable score with scale, the cost of capital intensity and the change since 2016--19; it then asks where the distance to the frontier sits. Descriptive statistics of reported scores use the baseline E1. Model-facing statistics are reported under E1, E2 and both imputation bounds (Section~\ref{sec:data-estimands}), and a claim is stated only where it holds in all four. Table~\ref{tab:estimands_v3} collects them.

\begin{table}[p]\centering\footnotesize\setlength{\tabcolsep}{1.5pt}\renewcommand{\arraystretch}{0.92}
\caption{Headline statistics on the audited sample (v3): E1 baseline, E2 and imputation bounds}\label{tab:estimands_v3}
\begin{tabular}{@{}>{\raggedright\arraybackslash}p{3.9cm}l>{\centering\arraybackslash}p{1.9cm}>{\centering\arraybackslash}p{1.9cm}>{\centering\arraybackslash}p{1.9cm}>{\centering\arraybackslash}p{1.9cm}c@{}}\toprule
 & & E1: all & E2: & \multicolumn{2}{c}{E2 plus imputed} & Holds in\\\cmidrule(lr){5-6}
Statistic & Basis & firm-years & organic & $M=4$ & $M=8$ & all four\\\midrule
\multicolumn{7}{l}{\emph{A. Model-facing statistics: E1 baseline; E2 and the imputation bounds (D-3)}}\\
$\Psi_{0.90}$ slope, 2010--25, $g$ 0--60\% (primary) & GAAP & $-$0.22 (0.11) & $-$0.29 (0.14) & $-$0.39 (0.11) & $-$0.36 (0.10) & yes/yes \\
 & cash & $-$0.16 (0.10) & $-$0.22 (0.11) & $-$0.25 (0.08) & $-$0.26 (0.08) & yes/no (3/4) \\
$\Psi_{0.90}$ slope, 2010--25, $g$ 0--40\% & GAAP & $-$0.31 (0.10) & $-$0.49 (0.14) & $-$0.44 (0.12) & $-$0.40 (0.12) & yes/yes \\
 & cash & $-$0.27 (0.10) & $-$0.37 (0.11) & $-$0.39 (0.09) & $-$0.37 (0.08) & yes/yes \\
$\Psi_{0.90}$ slope, 2010--19, $g$ 0--40\% & GAAP & $-$0.22 (0.10) & $-$0.28 (0.17) & $-$0.29 (0.13) & $-$0.28 (0.12) & yes/yes \\
 & cash & $-$0.16 (0.13) & $-$0.39 (0.18) & $-$0.37 (0.15) & $-$0.38 (0.13) & yes/yes \\
$\Psi_{0.90}$ slope, 2020--25, $g$ 0--40\% & GAAP & $-$0.60 (0.11) & $-$0.66 (0.12) & $-$0.60 (0.10) & $-$0.58 (0.11) & yes/yes \\
 & cash & $-$0.35 (0.08) & $-$0.34 (0.11) & $-$0.35 (0.08) & $-$0.33 (0.08) & yes/yes \\
$\Psi_{0.90}$ slope, 2020--25 minus 2010--19 & GAAP & $-$0.38 (0.13) & $-$0.37 (0.21) & $-$0.31 (0.16) & $-$0.30 (0.15) & yes/no (3/4) \\
 & cash & $-$0.19 (0.14) & 0.05 (0.20) & 0.02 (0.16) & 0.05 (0.15) & no (1/4)/no (0/4) \\
Local $Q_{0.90}$ slope, 2010--25 (descriptive) & GAAP & $-$0.41 (0.14) & $-$0.61 (0.17) & $-$0.51 (0.15) & $-$0.46 (0.14) & yes/yes \\
 & cash & $-$0.35 (0.11) & $-$0.43 (0.12) & $-$0.42 (0.09) & $-$0.39 (0.09) & yes/yes \\
Local $Q_{0.90}$ slope, 2010--19 & GAAP & $-$0.29 (0.14) & $-$0.43 (0.21) & $-$0.41 (0.17) & $-$0.34 (0.16) & yes/yes \\
 & cash & $-$0.20 (0.15) & $-$0.48 (0.19) & $-$0.48 (0.16) & $-$0.42 (0.16) & yes/yes \\
Local $Q_{0.90}$ slope, 2020--25 & GAAP & $-$0.71 (0.15) & $-$0.79 (0.15) & $-$0.68 (0.15) & $-$0.70 (0.14) & yes/yes \\
 & cash & $-$0.40 (0.10) & $-$0.42 (0.13) & $-$0.37 (0.09) & $-$0.41 (0.09) & yes/yes \\
Spend-to-excess-growth ratio (median) & S\&M & 3.02 (0.19) & 3.47 (0.25) & 3.34 (0.22) & 3.15 (0.23) &  \\
P95 of $R$ per tenfold revenue (points) & GAAP & 2.4 (2.8) & 4.0 (3.0) & 3.4 (2.9) & 3.2 (2.9) &  \\
 & cash & $-$4.1 (3.4) & $-$1.1 (3.9) & $-$2.5 (3.3) & $-$2.8 (3.1) & yes/no (0/4) \\
Capital-intensity gap in P95 of $R$, 2023--25 & GAAP & $-$1.8 [$-$9.5, +4.5] & $-$0.7 [$-$7.9, +5.4] & $-$0.8 [$-$8.3, +5.7] & $-$1.2 [$-$8.8, +5.3] &  \\
 & cash & $-$12.5 [$-$20.1, $-$3.9] & $-$7.4 [$-$18.3, $-$2.5] & $-$8.3 [$-$17.1, $-$2.6] & $-$12.3 [$-$17.9, $-$3.1] & yes/yes \\
 & capex added back & $-$4.6 [$-$14.6, +6.0] & $-$0.2 [$-$12.3, +9.0] & +0.1 [$-$10.6, +8.5] & $-$4.2 [$-$11.9, +7.0] & no (3/4)/no (0/4) \\
Change in $T(0\%)$, 2016--19 to 2023--25 & GAAP & +0.8 [$-$4.7, +4.5] & +0.4 [$-$4.9, +5.4] & +0.4 [$-$4.9, +4.9] & +0.8 [$-$4.8, +4.9] & yes \\
 & cash & $-$3.3 [$-$6.4, $-$0.2] & $-$2.5 [$-$6.5, +1.2] & $-$3.3 [$-$6.4, $-$0.4] & $-$3.2 [$-$6.2, +0.1] & no (2/4) \\
Change in $T(20\%)$, 2016--19 to 2023--25 & GAAP & $-$10.2 [$-$15.1, +1.1] & $-$9.7 [$-$15.1, +3.0] & $-$7.2 [$-$12.8, +4.2] & $-$6.3 [$-$12.0, +3.7] & yes \\
 & cash & $-$3.0 [$-$10.9, +3.4] & $-$1.1 [$-$7.7, +13.0] & $-$1.3 [$-$8.7, +8.6] & $-$1.8 [$-$8.7, +5.8] & yes \\
Change in $T(40\%)$, 2016--19 to 2023--25 & GAAP & $-$9.9 $^\dagger$ & -- & +1.0 $^\dagger$ & $-$2.0 $^\dagger$ &  \\
 & cash & +1.6 $^\dagger$ & -- & +7.6 $^\dagger$ & +7.2 $^\dagger$ &  \\
Near-frontier margin gap, average of bands & GAAP & 34.5 [28.6, 37.7] & 32.7 [26.6, 35.8] & 33.5 [28.0, 36.6] & 34.7 [28.9, 38.0] & yes/yes \\
\quad of which selling and marketing &  & 16.0 [12.4, 20.0] & 16.8 [12.6, 20.3] & 17.1 [12.3, 19.9] & 16.1 [12.5, 20.0] &  \\
\quad of which R\&D &  & 7.1 [4.6, 8.8] & 6.9 [4.0, 8.5] & 6.5 [4.2, 8.1] & 7.1 [4.6, 8.5] &  \\
\quad of which gross margin &  & 5.7 [$-$1.4, 9.4] & 2.9 [$-$1.7, 7.8] & 4.0 [$-$0.7, 8.7] & 4.9 [$-$0.7, 9.2] &  \\
\quad of which G\&A &  & 5.5 [3.4, 7.5] & 6.2 [3.3, 7.6] & 5.7 [3.4, 7.4] & 5.8 [3.4, 7.4] &  \\
\bottomrule\end{tabular}
\par\vspace{2pt}\footnotesize\raggedright Notes: Panels B and C and the notes follow on the next page. Cash: free cash flow after SBC.
\end{table}
\begin{table}[p]\centering\footnotesize\setlength{\tabcolsep}{1.5pt}\renewcommand{\arraystretch}{0.92}\addtocounter{table}{-1}
\caption{Headline statistics on the audited sample (v3), continued}
\begin{tabular}{@{}>{\raggedright\arraybackslash}p{3.9cm}l>{\centering\arraybackslash}p{1.9cm}>{\centering\arraybackslash}p{1.9cm}>{\centering\arraybackslash}p{1.9cm}>{\centering\arraybackslash}p{1.9cm}c@{}}\toprule
 & & E1: all & E2: & \multicolumn{2}{c}{E2 plus imputed} & Holds in\\\cmidrule(lr){5-6}
Statistic & Basis & firm-years & organic & $M=4$ & $M=8$ & all four\\\midrule
\multicolumn{7}{l}{\emph{A. Model-facing statistics (continued)}}\\
Anchor: median $\Pi$, all S\&M buys growth & cash & $-$1.3 [$-$18.2, +16.4] & $-$7.0 [$-$17.0, +15.4] & $-$0.5 [$-$13.6, +14.5] & $-$0.7 [$-$13.7, +15.9] & yes \\
Anchor: median $\Pi$, no S\&M buys growth & cash & +19.5 [+8.1, +26.5] & +18.5 [+6.8, +26.0] & +20.0 [+11.3, +26.9] & +20.2 [+12.2, +27.1] & yes/yes \\
Anchor: median $\Pi$, maintenance-netted (central) & cash & +12.4 [+0.3, +25.7] & +12.2 [+0.5, +24.7] & +13.2 [+3.5, +26.1] & +13.2 [+3.7, +26.5] & yes/yes \\
\addlinespace\multicolumn{7}{l}{\emph{B. Distribution of reported scores (descriptive, E1 primary)}}\\
P95 of $R$, 2016--19 & GAAP & 55.9 [51.1, 63.7] & 54.3 [46.6, 61.5] & 51.1 [46.9, 56.8] & 51.8 [47.8, 57.1] &  \\
 & cash & 57.3 [51.1, 67.3] & 53.9 [46.6, 60.6] & 51.7 [46.7, 57.3] & 53.1 [48.1, 57.4] &  \\
P95 of $R$, 2023--25 & GAAP & 43.9 [41.5, 47.2] & 42.8 [40.4, 45.1] & 43.2 [40.7, 45.7] & 43.5 [41.1, 46.0] &  \\
 & cash & 44.6 [41.5, 50.6] & 43.1 [39.4, 49.1] & 43.4 [40.2, 48.8] & 43.9 [40.6, 50.6] &  \\
Change in P95 of $R$, 2016--19 to 2023--25 & GAAP & $-$11.9 [$-$20.3, $-$6.4] & $-$11.4 [$-$18.7, $-$3.0] & $-$7.8 [$-$14.6, $-$2.7] & $-$8.3 [$-$15.1, $-$3.7] & yes/yes \\
 & cash & $-$12.7 [$-$21.6, $-$4.2] & $-$10.9 [$-$17.7, $-$2.1] & $-$8.3 [$-$15.2, $-$1.9] & $-$9.3 [$-$15.1, $-$1.3] & yes/yes \\
\addlinespace\multicolumn{7}{l}{\emph{C. Samples}}\\
Firm-years (firms) &  & 1,788 (292) & 1,378 (277) & 1,786 (291) & 1,786 (291) &  \\
\quad of which acquisition-affected &  & 466 & 58 & 466 & 466 &  \\
Firm-years with $g>50\%$ &  & 106 & 62 & 89 & 92 &  \\
\bottomrule\end{tabular}
\par\vspace{2pt}\footnotesize\raggedright Notes: Cash: free cash flow after SBC. The E1 $-$ E2 contrasts and the exclusion-only sample are in Table~\ref{tab:estimands_v3_more}; the maintenance-netted row nets the maintenance level of selling and marketing out of growth spend in each view and each draw (Section~\ref{sec:horizons-ranges}). Audited core sample (data/panel\_v3.csv: SIC 7372, fiscal years 2010--2025, firm-years first reported in their own annual report, captive related-party filer excluded, acquisition flag of rule R4). E1 uses every firm-year and reported growth (baseline). E2 keeps firm-years without material acquisitions and re-admits acquisition-affected ones with pro forma organic growth where the original filings identify it. The imputation columns add the acquisition-affected firm-years E2 drops, with organic growth imputed by removing acquired revenue equal to the purchase price over a price-to-revenue multiple $M$, half in the deal year and half in the next. $\Psi_{0.90}$ slope (primary statistic): least-squares slope over $g_0=0,0.05,\dots,0.6$ of $\phi_{0.90}(g_0)$, the 90th percentile of the margin among firm-years with $g\geq g_0$ (dominance cone; each grid point at least 30 firm-years in the sample and 10 in a bootstrap draw); within periods the cells at 60\% growth are below 30 firm-years, so period slopes use $0\leq g\leq0.4$. Local $Q_{0.90}$ slope (descriptive): $\tau=0.90$ linear quantile regression of margin on growth for $0\leq g\leq0.6$, solved exactly. Size gradient: $\tau=0.95$ quantile regression of $R$ on log prior-year revenue, times $\ln 10$. Capital-intensity gap: P95 of $R$ above minus below the period median of capital expenditure and capitalized software over revenue, split re-estimated in each bootstrap draw; firm-years without the measure are not classified. $T(g_0)=g_0+\phi_{0.95}(g_0)$. Anchor: 2023--25 firm-years at or above the P95 of the cash score (calibration rules). Parentheses: firm-cluster bootstrap standard errors; brackets: 90\% percentile intervals; 999 replications, firms drawn from E1, every column on the same draws; E1 $-$ E2 paired. $^\dagger$: no interval (more than 5\% of draws undefined, or a frontier-point cell below 30 firm-years). Last column (D-3): the claim's sign / its 90\% percentile interval hold in E1, E2 and both imputation columns; for frontier points at 0 and 20\% and the anchor's $\Pi_{ub}$ the claim is no change (interval includes zero). Points except slopes and the ratio. Bootstrap: 999 firm-cluster replications; one generator per block of statistics, seeded with 20261030 plus the CRC-32 checksum of the block name modulo $10^6$ (slopes 21088647, cone 20287800, spend 20456611, scale 20281130, capint 20868987, top 20806960, T 21225542, gap 20874387, anchor 20627171, sfa 20595418, grid 21165725, yearly 20451938, regimes 21094396, reweight 20547850, joint 20372013, bridge 20737029).
\end{table}


\subsection{Statistics}
Because both growth and margin are desirable, the relevant object is a Pareto (output--output) frontier of the points firms chose, which \autoref{def:industry} calls the observed frontier. Our primary statistic is its quantile version $\Psi_\chi$: for a growth level $g_0$, $\phi_\chi(g_0)$ is the $\chi$-quantile of the margin among firm-years that grew at least $g_0$ (computed over firm-years, so each firm contributes one chosen point a year), a partial frontier in the spirit of \citet{aragon2005nonparametric} and \citet{daraio2007advanced}, adapted to two desirable outputs. We summarize it by the least-squares slope of $\phi_{0.90}(g_0)$ over $g_0=0,0.05,\dots,0.60$. Each grid point needs at least \BootMinCell{} firm-years in the sample; only \PsiCellsEightyEone{} firm-years grow at least 80\%, so wider windows are not defined. Inside a bootstrap draw a grid point needs only 10 firm-years. Under E2, whose 60\% cell has \DefDrawOpCell{} firm-years in the sample, \DefDrawOpShare\% of the draws (GAAP) and \DefDrawCashShare\% (after SBC) evaluate that cell on fewer than the sample rule requires, more than the 5\% that would leave the statistic without an interval under that rule; on the draws that meet it, the E2 interval of the primary slope is \DefDrawOpLo{} to \DefDrawOpHi{} (GAAP) and \DefDrawCashLo{} to \DefDrawCashHi{} (after SBC), close to the published ones. E1 and both imputation bounds are unaffected. We report the slope for $\chi$ from 0.80 to 0.95 and windows of 0--40\% and 0--60\% (Table~\ref{tab:grid_v3}), and the order-$m$ frontier of \citet{cazals2002nonparametric}, the expected maximum margin of $m=25$ draws from the same set, as a check.

Alongside it we report the slope of a linear quantile regression \citep{koenker1978regression} of margin on growth at the 0.90 quantile for $0\leq g\leq0.6$, the local conditional quantile. It is a descriptive statistic, not the object of the theory. It compares margins at each growth rate, while $\Psi_\chi$ pools all faster growers, so it falls wherever margins fall with growth, and $\Psi_\chi$ only where the best margins of fast growers fall short of the best margins of slower ones. The best score attainable on the frontier, $\max_g[g+\phi(g)]$, would be a natural summary, but it sits at the edge of the growth support, where a handful of firm-years determine it, so we do not use it. We also report upper quantiles of the score $R=g+m$, frontier points at fixed growth thresholds, $T(g_0)=g_0+\phi_{0.95}(g_0)$, fixed-effects regressions and a normal--half-normal stochastic frontier \citep{aigner1977formulation,meeusen1977efficiency}.

Quantile regressions are solved exactly as linear programs. Intervals come from a bootstrap that resamples firms, preserving each firm's time series, with \BootB{} replications. Firms are drawn from the \BootFirms{} firms of E1 and every estimand is evaluated on its own firm-years of the drawn firms, so differences between estimands are paired. We report percentile intervals, and basic, normal and bias-corrected intervals where they differ. A statistic with more than \BootNanMax\% of undefined draws gets no interval. Familywise error is controlled by the step-down procedure of \citet{holm1979simple}, with the step-down max-$t$ procedure of \citet{romano2005stepwise} as a check, its statistics scaled by their bootstrap standard errors, in two families: the \ContrN{} differences between E1 and E2, and the \FamBN{} headline statistics of this section. Two members of the second family test the primary slopes against $-1$, a null the model does not hold, and the average margin gap reported below is positive by construction, so the survival of these three of the correction says little; the family was chosen after the first version's results were known. The survival statements depend on the family: adding the frontier-point changes $T(20\%)$ on both bases and $T(0)$ after SBC, and the anchor test of Section~\ref{sec:empirics-anchor}, gives an \FamEighteenN-test family in which the Holm-adjusted $p$-value of the 2020--25 steepening is \HolmEighteenDiffNormal, \HolmEighteenDiffPercentile{} and \HolmEighteenDiffBasic{} (normal, percentile, basic), that of the GAAP decline \HolmEighteenTopNormal, \HolmEighteenTopPercentile{} and \HolmEighteenTopBasic, and that of the anchor \HolmEighteenAnchorNormal, \HolmEighteenAnchorPercentile{} and \HolmEighteenAnchorBasic.

\subsection{The top of the distribution fell}
This subsection describes observed scores (E1). The median GAAP score fell from \VMedOpEarly{} in 2010--15 to \VMedOpLate{} in 2023--25, and the share of firm-years clearing 40 from \VPassOpEarly\% to \VPassOpLate\%; after deducting SBC from free cash flow the median fell from \VMedCashEarly{} to \VMedCashLate{} and the pass rate from \VPassCashEarly\% to \VPassCashLate\%.

The top of the distribution fell as well. Between 2016--19 and 2023--25 the 95th percentile of the GAAP score fell from \TopEarlyOp{} to \TopLateOp, by \TopFallOpEone{} points (90\% interval \TopFallOpEoneLo{} to \TopFallOpEoneHi), and after SBC from \TopEarlyCash{} to \TopLateCash, by \TopFallCashEone{} points (\TopFallCashEoneLo{} to \TopFallCashEoneHi). The decline holds in \TopCount{} views: on the GAAP basis it is \TopFallOpEtwo{} points under E2 and \TopFallOpImpFour{} and \TopFallOpImpEight{} under the imputation bounds. Against the family of headline statistics, the Holm-adjusted $p$-value of the GAAP decline is \HolmTopLo{} to \HolmTopHi{} across interval types (Romano--Wolf \HolmTopRW, with a Monte Carlo standard error of about \HolmTopRWSe{} at \BootB{} draws), so it survives at a familywise level of 10\% under the normal and percentile types and under Romano and Wolf's procedure, but not under the basic type, and at 5\% only under the normal type. The cash-basis decline survives the correction under no procedure, even at 10\%, although each of its unadjusted intervals excludes zero. The yearly series (Table~\ref{tab:yearly_v3}) shows that the decline is not a steady trend. The GAAP 95th percentile reached its low of \TopYearMinOp{} in 2024 and recovered to \TopYearLastOp{} in 2025. The recovery is a feature of observed scores: against 2023--24, the 2025 GAAP 95th percentile is \RecoveryEone{} points higher under E1 but only \RecoveryNoAcq{} without the acquisition-affected firm-years, which include several of the highest 2025 scores, and it is smaller under E2 (Table~\ref{tab:yearly_v3}). Measured from 2016--19 to 2023--24 the change is \TopTwentyFourOp{} points (\TopTwentyFourOpLo{} to \TopTwentyFourOpHi), and to 2025 alone \TopTwentyFiveOp{} (\TopTwentyFiveOpLo{} to \TopTwentyFiveOpHi).

The decline came with a change in the growth mix, but the data do not say how much of it the mix explains, or why it occurred. Only \FortyCellLate{} firm-years grew at least 40\% in 2023--25, against \FortyCellEarly{} in 2016--19, and \TwentyCellLate{} grew at least 20\%, against \TwentyCellEarly. A reweighting decomposition over ten growth bins, in the manner of \citet{dinardo1996labor}, splits the change into the two sources. Reweighting the 2016--19 firm-years to the growth mix of 2023--25 attributes \RwOpCompositionB{} points of the GAAP change to composition (\RwOpCompositionBLo{} to \RwOpCompositionBHi) and \RwOpWithinB{} to changes within growth bins (\RwOpWithinBLo{} to \RwOpWithinBHi); reweighting in the other direction gives \RwOpCompositionA{} and \RwOpWithinA. On the cash basis the composition effect is \RwCashCompositionB{} and \RwCashCompositionA{} in the two directions, and the within-bin change \RwCashWithinB{} and \RwCashWithinA. Reweighting to the 2016--19 mix of firm sizes instead leaves the decline in place (\RwOpWithinSize{} points on the GAAP basis), so the size mix does not account for it. The split is not robust: with 4 to 40 growth bins in either direction the composition part runs from \BinCompOpLo{} to \BinCompOpHi{} points on the GAAP basis and the within-bin part from \BinWithinOpLo{} to \BinWithinOpHi{} (after SBC \BinCompCashLo{} to \BinCompCashHi{} and \BinWithinCashLo{} to \BinWithinCashHi), the bootstrap cannot separate the two parts, and because the score contains growth, a reweighting on growth bins counts shifts of growth within the top bin as within-bin change. The fall was not confined to the top: on the GAAP basis the 10th percentile of the score changed by \QShiftOpTen{} points, the median by \QShiftOpFifty{} and the 95th percentile by \QShiftOpNinetyFive; on free cash flow before SBC the changes were \QShiftFcfTen, \QShiftFcfFifty{} and \QShiftFcfNinetyFive. Competing readings fit the same facts: the normalization of growth after 2020--22, the rise in the cost of capital, which in the model lowers a firm's own weight on growth (\autoref{prop:dynamic}, under (H)), and the closing of the window for new listings, which cut the firm-years in a firm's first year in the panel from \FirstYearTwentyTwentyTwo{} in 2020--22 to \FirstYearTwentyThreeTwentyFive{} in 2023--25. The data cannot separate these readings.

Frontier points at fixed growth thresholds do not show that the best margins at a given growth rate held up (Table~\ref{tab:estimands_v3}). The 95th-percentile GAAP margin among firm-years with nonnegative growth, $T(0)$, changed by \TzeroOp{} points between 2016--19 and 2023--25 (\TzeroOpLo{} to \TzeroOpHi), but among firm-years growing at least 20\%, $T(20\%)$, it changed by \TtwentyOp{}, about as much as the score itself; under E1 its percentile interval runs from \TtwentyOpLo{} to \TtwentyOpHi, while the basic, normal and bias-corrected intervals exclude zero. After SBC, $T(0)$ changed by \TzeroCash{} points (\TzeroCashLo{} to \TzeroCashHi), a fall whose interval excludes zero under E1 and in two of the four views, and $T(20\%)$ changed by \TtwentyCash{} (\TtwentyCashLo{} to \TtwentyCashHi). The data therefore cannot distinguish a change in the growth mix from a fall in the best margins of firms growing 20\% or more. We do not report the frontier point at 40\%. Its 2023--25 cell has \FortyCellLate{} firm-years under E1 and \FortyCellLateEtwo{} under E2, and \FortyNanEone\% of E1 draws and \FortyNanEtwo\% of E2 draws are undefined, above the limit of \BootNanMax\%; the first version's finding that the organic 40\% point rose is withdrawn.

\subsection{At the top, a point of growth costs well under a point of margin}
The primary slope is that of $\Psi_{0.90}$ over growth of 0 to 60\%. It is \PsiOpEone{} (standard error \PsiOpEoneSe; 90\% interval \PsiOpEoneLo{} to \PsiOpEoneHi) on the GAAP basis and \PsiCashEone{} (\PsiCashEoneSe; \PsiCashEoneLo{} to \PsiCashEoneHi) after deducting SBC from free cash flow. Under E2 it is \PsiOpEtwo{} and \PsiCashEtwo, and under the imputation bounds \PsiOpImpFour{} and \PsiOpImpEight{} on the GAAP basis and \PsiCashImpFour{} and \PsiCashImpEight{} after SBC. Across $\chi$ from 0.80 to 0.95 and both windows the E1 slope lies between \PsiGridOpLo{} and \PsiGridOpHi{} on the GAAP basis and between \PsiGridCashLo{} and \PsiGridCashHi{} after SBC, and across all four views and both bases between \PsiGridAllOpLo{} and \PsiGridAllCashHi; the order-$m$ slopes are \OrderMOp{} and \OrderMCash.

Two claims follow, with different strength. That the slope is far above $-1$, the one-for-one trade that the Rule-of-40 line describes, holds in \PsiOpCountAbove{} views and survives Holm's correction at a familywise level of 5\% on both bases. That it is negative holds in sign in \PsiOpSignNeg{} views, and its 90\% interval excludes zero in \PsiOpCountNeg{} views on the GAAP basis and \PsiCashCountNeg{} after SBC; but it does not survive the familywise correction (Holm-adjusted $p$-values \HolmPsiOpZeroLo{} to \HolmPsiOpZeroHi{} on the GAAP basis and \HolmPsiCashZeroLo{} to \HolmPsiCashZeroHi{} after SBC). Among the best-performing software firms, then, a point of growth has come with well under a point less margin, about a fifth of a point in the baseline. The local conditional quantile, reported as a descriptive statistic, is steeper: \LocOpEone{} (\LocOpEoneSe) on the GAAP basis and \LocCashEone{} (\LocCashEoneSe) after SBC under E1, and \LocOpEtwo{} and \LocCashEtwo{} under E2. It is steeper here because it compares margins at a given growth rate, where the typical fast grower is less profitable than the typical slow one, while the frontier compares the best of each.

Did the trade-off steepen? Within periods the cells at 60\% growth are too thin, so we compare periods on the 0--40\% window. On the GAAP basis the $\Psi_{0.90}$ slope was \PsiEarlyOp{} (\PsiEarlyOpSe) in 2010--19 and \PsiLateOp{} (\PsiLateOpSe) in 2020--25, a difference of \PsiDiffOpEone{} (\PsiDiffOpEoneSe; \PsiDiffOpEoneLo{} to \PsiDiffOpEoneHi). The difference is negative in \PsiDiffOpSignNeg{} views, and its 90\% interval excludes zero in \PsiDiffOpCountNeg{} of them. Its Holm-adjusted $p$-value is \HolmPsiDiffLo{} to \HolmPsiDiffHi{} (Romano--Wolf \HolmPsiDiffRW), so it survives at a familywise level of 10\% under every procedure and at 5\% under \HolmDiffFive{} of the four. After SBC there was no steepening: the difference is \PsiDiffCashEone{} (\PsiDiffCashEoneSe) under E1 and positive in the other three views. The GAAP frontier steepened after 2020 in sign and with an interval that excludes zero in \PsiDiffOpCountNeg{} views; the cash frontier did not. In 2025 the within-year slope of the local quantile, steep in 2020--24, was close to zero (Table~\ref{tab:yearly_v3}).

The cross-sectional slope is far flatter than the slope of any firm's own frontier at its chosen point, which by \autoref{prop:optimum} equals $-X$, and $X\geq1$ for a firm that values a point of growth at least as much as a point of margin. \autoref{prop:industry} explains why the two can differ: every firm faces a downward-sloping frontier, but the estimated frontier is drawn through the points firms chose, not through their possibilities, and when firms with more growth opportunity also have higher margin intercepts the chosen points slope down much less than any individual frontier. The cross-section does not identify any firm's marginal cost of growth (\autoref{prop:industry}(iv)); a common frontier with high opportunity can produce a similar slope. Within firms, margin and growth are nearly unrelated: with firm and year fixed effects, standard errors clustered by firm and growth up to 150\%, the coefficient of GAAP margin on growth is \WfOpNow{} (\WfOpNowSe), and after SBC \WfCashNow{} (\WfCashNowSe). In the model, spend this year buys growth next year, so we also relate this year's margin to next year's growth: the coefficients are \WfOpLead{} (\WfOpLeadSe) and \WfCashLead{} (\WfCashLeadSe) under E1, and \WfEtwoOpLead{} (\WfEtwoOpLeadSe) and \WfEtwoCashLead{} (\WfEtwoCashLeadSe) under E2 when both years belong to the organic sample, and \WfImpFourOpLeadSame{} and \WfImpEightOpLeadSame{} (GAAP) and \WfImpFourCashLeadSame{} and \WfImpEightCashLeadSame{} (after SBC) under the imputation bounds. These are small. \autoref{prop:industry}(v) shows that anticipated changes in opportunity would produce a strongly negative within-firm slope, so the weak relation suggests that growth surprises against costs already committed, or demand shocks that also raise margins, offset much of the cost of buying growth; the regressions do not identify which.

\subsection{What the data do and do not say about the growth technology}
The model's growth technology has two parameters, demand pull and the spend elasticity $\alpha$, and the data do not identify them separately. Under \autoref{prop:optimum}, an optimizing firm spends $\sigma^*=\alpha(X-k)\Delta^*$, so the ratio of spend to excess growth reveals at most the product $\alpha(X-k)$, and only if spend and excess growth are measured as the model defines them. With selling and marketing expense as the measure of spend, the median ratio of the previous year's selling-and-marketing share, taken like growth from the same filing, to growth in excess of the year's median, among firm-years growing faster than the median, is \SpendEone{} (\SpendEoneSe) under E1, \SpendEtwo{} under E2 and \SpendImpFour{} and \SpendImpEight{} under the imputation bounds. This spend ratio is robust across the four views, at \SpendLo{} to \SpendHi, and we report it as a statistic of its own.

It does not pin down $\alpha(X-k)$, because the answer depends on how excess growth and the spend that buys it are measured. By terciles of excess growth the ratio is \AlphaTercLow, \AlphaTercMid{} and \AlphaTercHigh; among firm-years growing more than 10 or 20 points above the median it is \AlphaFloorTen{} and \AlphaFloorTwenty; netting out the median selling-and-marketing share of firm-years growing at or below the median, as maintenance spend, gives \AlphaNet; regressions of the selling-and-marketing share on growth, which use the variation in levels, give \AlphaLevelLo{} to \AlphaLevelHi; and replacing the year's median growth by a fixed market growth of 5, 8 or 12\% gives \AlphaGFive, \AlphaGEight{} and \AlphaGTwelve. Across methods $\alpha(X-k)$ ranges from \AlphaAllLo{} to \AlphaAllHi{} under E1, so for $\alpha$ between 0.25 and 0.45 the implied $X-k$ ranges from about \XkLo{} to about \XkHi. The data therefore do not identify the weight on growth that firms act on, and we withdraw the first version's statement that it lies between 7 and 13. What remains is the theoretical distinction. Under (H), the weight a value-maximizing firm puts on growth is the marginal value of a dollar of revenue (\autoref{prop:dynamic}), a valuation multiple, while an investor's screen applies a duration-based weight of two to three; the two need not coincide, and the cross-section of chosen points cannot reveal either. We treat $\alpha$ as a calibration parameter in $[0.25,0.45]$ and anchor the calibration of Section~\ref{sec:horizons} on observed firms rather than on a regression.

\subsection{Scale}
\autoref{prop:saturation} predicts that the maximal score falls as market share rises. The evidence is weak and not robust. After SBC, the 95th percentile of the score is \ScaleClassA{} for firms with \$0.1--0.5 billion of prior revenue, \ScaleClassB{} and \ScaleClassC{} in the next two classes and \ScaleClassD{} above \$10 billion. A 0.95-quantile regression on log prior revenue gives \ScaleCashEone{} points per tenfold increase (standard error \ScaleCashEoneSe) under E1, \ScaleCashEtwo{} (\ScaleCashEtwoSe) under E2 and \ScaleCashImpFour{} and \ScaleCashImpEight{} under the imputation bounds; with period effects the E1 and E2 gradients are \ScaleCashFeEone{} and \ScaleCashFeEtwo. The 90\% interval includes zero in every view, with and without period effects, and on the GAAP basis the E1 gradient is \ScaleOpEone. Size is in any case a noisy proxy for market share, and larger firms enjoy operating leverage that the model omits, so the data do not test saturation.

\subsection{Capital intensity: an accounting effect}
\autoref{prop:capital} implies that on a cash basis capital intensity lowers the maximal score in proportion to growth, while operating margins are affected only through depreciation. We split each period at its median ratio of capital expenditure and capitalized software to revenue, re-estimating the split in every bootstrap draw, and compare the 95th percentile of the score. In 2023--25 the difference between the capital-intensive half and the rest is \CapCashEone{} points after SBC (\CapHiHalf{} against \CapLoHalf; \CapCashEoneLo{} to \CapCashEoneHi), an interval that excludes zero in \CapCashCount{} views, but only \CapOpEone{} points on the GAAP basis (\CapOpEoneLo{} to \CapOpEoneHi). The gap is mostly mechanical. The sorting variable is deducted from the cash margin, and adding it back leaves \CapAddEone{} points (\CapAddEoneLo{} to \CapAddEoneHi) under E1 and \CapAddEtwo{} under E2; sorting on each firm's median intensity, which a single year's spending cannot move, gives \CapFirmMedEone{} (\CapFirmMedEoneLo{} to \CapFirmMedEoneHi). In 2016--19 the add-back removes a smaller part of the gap (\CapCashEarlyEone{} points, and \CapAddEarlyEone{} with the deduction added back, \CapAddEarlyEoneLo{} to \CapAddEarlyEoneHi), and no sort or basis gives an interval that excludes zero in all four views. The prediction that the gap grows with growth is not supported either: in a 0.90-quantile regression of the cash margin, the interaction of growth with high intensity is \CapInter{} (\CapInterSe), negative in only \CapInterSign{} of the four views. We therefore report the cash-basis gap as an accounting effect of deducting capital spending, not as evidence for \autoref{prop:capital}. CoreWeave, described in Section~\ref{sec:data}, illustrates how large that deduction can be.

\subsection{The distance to the frontier}
The stochastic frontier measures how far the typical firm sits below the frontier. We estimate one frontier per period, common to all cohorts, and the conditional mean of each firm-year's shortfall given its residual \citep{jondrow1982estimation}. The mean shortfall on the GAAP basis was \SfaEarly{} margin points in 2016--19 and \SfaLate{} in 2023--25. In 2023--25, the shortfall of firms that entered the sample in 2019 or later exceeds that of earlier cohorts by \SfaCohortOp{} points on the GAAP basis (\SfaCohortOpLo{} to \SfaCohortOpHi) and \SfaCohortCash{} after SBC (\SfaCohortCashLo{} to \SfaCohortCashHi), with an interval excluding zero in \SfaCohortCount{} views. Firms present in both periods did not fall further behind: their shortfall changed by \SfaBalanced{} points (\SfaBalancedLo{} to \SfaBalancedHi). Two caveats apply. Only \SfaStill{} of the \SfaPreFirms{} firms that entered before 2019 and appear in 2016--19 remain in 2023--25, and their 2016--19 median GAAP margin was \SfaSurvMargin\%, against \SfaLeaveMargin\% for those that left, so the comparison sets survivors against entrants whose exits are not yet observed. And the shortfall is a distance below a parametric frontier at given growth, not a measure of managerial inefficiency.

Because operating margin is gross margin minus selling and marketing, research and development, general and administrative and other operating expense, the margin gap between firms near the frontier and typical firms at the same growth rate splits exactly into gaps in each line. Within each growth band we compare firm-years in the top decile of GAAP operating margin with those between the 40th and 60th percentiles, for the firm-years with growth of 0 to 60\% that report all four lines. On average across bands the top decile earns \GapV{} points more margin [\GapVLo, \GapVHi] under E1 (\GapVEtwo{} under E2). Selling and marketing accounts for \GapVSM{} points [\GapVSMLo, \GapVSMHi], research and development for \GapVRD{} and general and administrative expense for \GapVGA; the gross-margin part, \GapVGM{} [\GapVGMLo, \GapVGMHi], is not robust, with a percentile interval that includes zero in every view (under E1 the basic, normal and bias-corrected intervals exclude zero). The composition changes with growth. Among firms growing less than 10\% a year, gross margin accounts for \GapVSlowGM{} of \GapVSlow{} points. Among firms growing 40 to 60\%, selling and marketing accounts for \GapVFastSM{} of \GapVFast{} points [\GapVFastSMLo, \GapVFastSMHi] and research and development for another \GapVFastRD: the best fast growers buy the same growth with far less selling and marketing expense, which in the model is higher demand pull. The average gap and the selling-and-marketing part among fast growers survive Holm's correction at a familywise level of 5\%, but the first is positive by construction and the second is a term of an accounting identity, so the correction adds little to the description. The decomposition is descriptive. Firms near the frontier may spend less on selling because their products pull demand, not because they sell more efficiently, so the decomposition shows where margins differ, not which changes would close the gap. It agrees, however, with the ordering of \autoref{cor:priorities}: the growth engine matters most for fast growers and overhead for slow growers, while the corollary values savings in cost of revenue only to the extent pricing lets the firm keep them.

\subsection{The calibration anchor}\label{sec:empirics-anchor}
Section~\ref{sec:horizons} anchors the calibration on the firm-years at or above the 95th percentile of the cash score in 2023--25, without acquisition-affected years: \AncNEone{} firm-years of \AncFirmsEone{} firms under E1 and \AncNEtwo{} of \AncFirmsEtwo{} under E2. For each it computes the growth surplus $\Pi$, the excess of the score over its saturated level, under two bounds on how much selling spend buys growth. When none does, the median surplus is \AncPiLbEone{} points under E1 [\AncPiLbEoneLo, \AncPiLbEoneHi] and \AncPiLbEtwo{} under E2 [\AncPiLbEtwoLo, \AncPiLbEtwoHi], with an interval above zero in \AncPiLbCount{} views. When all of it does, the median is \AncPiUbEone{} [\AncPiUbEoneLo, \AncPiUbEoneHi] under E1 and \AncPiUbEtwo{} [\AncPiUbEtwoLo, \AncPiUbEtwoHi] under E2, an interval that includes zero in \AncPiUbCount{} views; leaving out one firm at a time moves the E1 median between \AncLooLo{} and \AncLooHi. On the central basis, which nets the maintenance level of selling and marketing out of growth spend (Section~\ref{sec:horizons-ranges}), the median is \AncPiNetEone{} [\AncPiNetEoneLo, \AncPiNetEoneHi] under E1 and \AncPiNetEtwo{} [\AncPiNetEtwoLo, \AncPiNetEtwoHi] under E2, above zero in \AncPiNetCount{} views but with a lower end near zero under E1 and no adjustment for multiplicity (Table~\ref{tab:estimands_v3}). Only the statement in which no selling spend buys growth is a conclusion: if selling spend does not buy growth, the best scores of 2023--25 exceed their durable level. It lies outside both familywise families; added to the family of headline statistics, its Holm-adjusted $p$-value over \FamBNAnchor{} tests (the family, with its non-estimable member at $p=1$ as in the \FamEighteenN-test family above, and the anchor) is \HolmAnchorFifteenNormal{} (normal), \HolmAnchorFifteenPercentile{} (percentile) and \HolmAnchorFifteenBasic{} (basic), so it survives at 10\% under every type and at 5\% under two; in the larger family above it does not survive at 10\% under the percentile type (\HolmEighteenAnchorPercentile). Two firms supply six of the \AncNEone{} firm-years; Table~\ref{tab:anchor-firms} of the Supplementary Appendix names them. The share of the median top score that is transient is \AncTrUbEone\% [\AncTrUbEoneLo, \AncTrUbEoneHi] at the upper bound on spend and \AncTrLbEone\% [\AncTrLbEoneLo, \AncTrLbEoneHi] at the lower bound under E1, and between \AncTrAllLo\% and \AncTrAllHi\% across bounds and views, so the data do not determine it.

\subsection{Robustness}\label{sec:empirics-robust}
Table~\ref{tab:robust} of the Supplementary Appendix repeats the headline statistics in \NVariants{} samples: the four views; the organic sample at acquisition flags of 10\% and 50\% of prior-year revenue; exclusion of every acquisition-affected year, at each flag; the two readings of ambiguous pro forma disclosures; the rule of the first test run; the excluded fiscal-2010 years added back; the samples without the firm-years flagged under rule R5 and without the debt-only filers a referee named; and the samples without the firms that both, or either, of the two independent readers of rule R2 class as not software and without all firms the review panels named. In every variant the 95th percentile falls by \TopFallAllLo{} to \TopFallAllHi{} points between 2016--19 and 2023--25, and the $\Psi_{0.90}$ slope lies between \PsiVarOpLo{} and \PsiVarOpHi{} on the GAAP basis and between \PsiVarCashLo{} and \PsiVarCashHi{} after SBC. Among the observed-score samples, dropping the firms that the readers of rule R2 class as not software moves the statistics most: without the firms both readers so class, the 95th percentile falls by \VRtwoBothFallOp{} points on the GAAP basis and \VRtwoBothFallCash{} after SBC, and the $\Psi_{0.90}$ slope is \VRtwoBothPsiOp{} and \VRtwoBothPsiCash. Table~\ref{tab:grid_v3} gives the slopes by quantile and window.

The estimands are shown side by side, but we make no claim that they differ (Table~\ref{tab:contrasts_v3}). Of the \ContrN{} paired differences between E1 and E2, \ContrExclPct{} have a 90\% percentile interval that excludes zero, \ContrExclBasic{} a basic interval and \ContrExclNormal{} a normal interval, against \ContrExpected{} expected by chance if no difference were real. None survives Holm's correction at a familywise level of 5\% under any interval type, nor Romano and Wolf's procedure; at 10\%, only the difference in the local GAAP slope survives, and only with basic $p$-values (adjusted $p$ \ContrLocalOpHolm). The first version's statements that acquired growth flattens the GAAP slope and that the size gradient and the steepening of the cash slope reflect acquired growth are withdrawn: imputing the dropped years returns values close to E1's, so the differences reflect which years E2 drops.

\subsection{Summary}
Six facts carry into the rest of the paper; each holds under observed scores, the organic sample and both imputation bounds unless stated otherwise. (i) Among the best-performing software firms a point of growth has come with well under a point less margin: the slope of the 0.90-quantile frontier is \PsiOpEone{} on the GAAP basis and \PsiCashEone{} after SBC in the baseline, far above $-1$ in every view; its negative sign holds in every view but does not survive the familywise correction. Because the cross-section records choices, the slope does not measure any firm's own trade-off. The GAAP slope steepened after 2020, in sign and with an interval that excludes zero in \PsiDiffOpCountNeg{} views; the cash-basis slope did not. (ii) The data identify neither the curvature of the growth technology nor the weight on growth that firms act on: the spend ratio of \SpendLo{} to \SpendHi{} is robust, but the product $\alpha(X-k)$ it is meant to reveal ranges from \AlphaAllLo{} to \AlphaAllHi{} across methods. (iii) A decline of the upper tail with scale is weak and not robust. (iv) The lower cash-basis scores of capital-intensive firms are an accounting effect of deducting capital spending. (v) Between 2016--19 and 2023--25 the 95th percentile of the score fell by \TopFallOpEone{} points on the GAAP basis and \TopFallCashEone{} after SBC, while fewer firms grew fast; the GAAP decline survives a familywise correction at 10\% under \TopFamTen{} of four procedures, and at 5\% under \TopFamFive, and the cash decline under none. The 95th-percentile GAAP margin of firms growing at least 20\% fell by about as much, and after SBC that of all firms with nonnegative growth fell in two of four views, so the data do not separate a change in the growth mix from a fall in the best margins of fast growers. The top partly recovered in 2025 under observed scores. (vi) Firms near the frontier earn about \GapV{} points more operating margin than typical firms at the same growth rate, mostly through lower selling, research and overhead costs, and recent entrants sit further below the frontier than earlier cohorts. Whether the best scores of 2023--25 are above their durable level depends on how much selling spend buys growth; the data say only that they are if little of it does.
